% ===================================================================
% FST-III: Biologische Spieltheorie -- Von molekularer ESS zu Verhaltensattraktoren
% Zieljournal: Journal of Theoretical Biology / BioSystems
% Version: 2.0 (integriertes Quellmaterial)
% ===================================================================

\documentclass[12pt,a4paper]{article}

% Packages
\usepackage[utf8]{inputenc}
\usepackage[T1]{fontenc}
\usepackage{amsmath,amssymb,amsthm}
\usepackage{physics}
\usepackage{graphicx}
\usepackage{hyperref}
\usepackage{booktabs}
\usepackage[margin=2.5cm]{geometry}
\usepackage{natbib}
\usepackage{cleveref}
\usepackage{enumitem}

% Theorem environments
\newtheorem{proposition}{Proposition}
\newtheorem{conjecture}{Conjecture}
\newtheorem{definition}{Definition}
\theoremstyle{remark}
\newtheorem{remark}{Remark}

% Title
\title{Biologische Spieltheorie:\\
Von molekularer ESS zu Verhaltensattraktoren ---\\
Ein vereinigender Rahmen zur Verkn\"upfung von Thermodynamik und evolution\"arer Spieltheorie}

\author{Lukas Geiger\\
\textit{Unabh\"angiger Forscher, Bernau im Schwarzwald, Deutschland}\\
ORCID: \url{https://orcid.org/0009-0005-7296-1534}}

\date{M\"arz 2026}

\begin{document}

\maketitle

\begin{abstract}
Dieses Paper pr\"asentiert einen vereinigenden Rahmen f\"ur biologische Organisation von der molekularen bis zur \"Okosystem-Skala, begr\"undet in thermodynamischer Spieltheorie und dem Maximum Entropy Production Principle (MEPP). Wir argumentieren, dass biologische Systeme hocheffiziente dissipative Strukturen sind, deren stabile Konfigurationen als Nash-Gleichgewichte auf multiplen Organisationsebenen charakterisiert werden k\"onnen. Der Rahmen integriert Jeremy Englands Theorie der dissipativen Adaptation, Karl Fristons Free Energy Principle, Stuart Kauffmans Boolesche Netzwerkdynamik und die klassische evolution\"are Spieltheorie in eine koh\"arente Hierarchie. Auf jeder Ebene---von der Proteinfaltung \"uber Genregulation, zellul\"are Koalitionen, organismisches Verhalten bis zur \"Okosystemdynamik---regeln Nash-Gleichgewichte stabile Konfigurationen. Die Haupt\"uberg\"ange der Evolution, identifiziert von Szathm\'ary und Maynard Smith, werden als Grobk\"ornungsschritte uminterpretiert, die spieltheoretische Struktur \"uber Skalen hinweg erhalten. Wir f\"uhren das Konzept der \emph{Nash-Frustration} als Komplement zur energetischen Frustration in Proteinstrukturen ein und leiten testbare Vorhersagen f\"ur Proteinfaltungslandschaften, Genregulationsnetzwerk-Attraktoren und \"Okosystem-Resilienz ab. Spekulative kosmologische Erweiterungen werden separat in Abschnitt~\ref{sec:cosmology} diskutiert und sollten als vorl\"aufige Forschungsagenda gelesen werden, nicht als etablierte Theorie.

\textbf{Hinweis zum G\"ultigkeitsbereich:} Die prim\"aren Beitr\"age dieses Papers betreffen biologische Skalen (Abschnitte~\ref{sec:thermodynamics}--\ref{sec:ecosystem}). Die kosmologische Diskussion (Abschnitt~\ref{sec:cosmology}) ist explizit spekulativ und behandelt Verbindungen, die eine dedizierte Behandlung in einer separaten Arbeit erfordern w\"urden.
\end{abstract}

\textbf{Schl\"usselw\"orter:} evolution\"are Spieltheorie, ESS, Proteinfaltung, Genregulation, MEPP, dissipative Adaptation, Free Energy Principle, Boolesche Netzwerke, Haupt\"uberg\"ange, Renormierungsgruppe

% ===================================================================
\section{Einleitung}
\label{sec:introduction}
% ===================================================================

Die Entstehung und Stabilisierung biologischer Komplexit\"at stellt die moderne Wissenschaft vor die Herausforderung, die fundamentalen Gesetze der Physik und die dynamischen Strategien des Lebens zu verbinden. Der Rahmen der Functional Stability Theory (FST-III) postuliert, dass biologische Systeme nicht blo\ss{} Produkte zuf\"alliger Mutationen sind, sondern hocheffiziente dissipative Strukturen, die innerhalb eines strikten thermodynamischen und kosmologischen Rahmens operieren \citep{EcosystemMEPP}.

Dieses Paper untersucht die vorgeschlagene Verbindung zwischen Spieltheorie, dem Maximum Entropy Production Principle (MEPP) und kosmologischen Randbedingungen und erforscht, ob Leben als Attraktor innerhalb physikalischer Einschr\"ankungen verstanden werden kann.

Die klassische evolution\"are Spieltheorie, bahnbrechend von Maynard Smith und Price \citep{MaynardSmith1982}, revolutionierte unser Verst\"andnis biologischen Verhaltens. Doch die Standardbehandlung beschr\"ankt sich weitgehend auf die Verhaltens\"okologie. Dieses Paper argumentiert f\"ur eine breitere Vision: spieltheoretische Prinzipien operieren auf jeder Ebene biologischer Organisation, und diese Prinzipien verbinden die Biologie mit der Physik durch die gleiche Optimierungslogik, die die Teilchenphysik \citep{FSTI} und die pr\"abiotische Chemie \citep{FSTII} regiert.

% ===================================================================
\section{Strukturelle Ausrichtung: FST-III als Pattern-A-Instanz}
\label{sec:structural-alignment}
% ===================================================================

Der hier entwickelte biologische Stabilit\"atsrahmen (ESS/FEP) ist nicht blo\ss{} eine Heuristik, sondern eine Manifestation der \textbf{Pattern-A-Stabilit\"atslogik}, die das FST-\"Okosystem vereinigt (siehe \cite{FST-Unified2026}). Die j\"ungste \textbf{eichtheoretische Reformulierung der Riemannschen Vermutung} (siehe \cite{FST-RH3}) liefert das strukturelle Vorbild f\"ur diese Architektur.

Innerhalb dieses hierarchischen Stabilit\"atsrahmens werden biologische Systeme wie folgt verstanden:
\begin{enumerate}[label=\alph*)]
    \item \textbf{Analytische Rang-Endlichkeit (Null-Tail):} Biologische Komplexit\"at ist gegen die kombinatorische Explosion des ``toten'' Zustands gesch\"utzt, weil nur eine endliche Teilmenge von Genotypen ESS/FEP-Stabilit\"at erreicht. Das unorganisierte ``Rauschen'' wird durch die entropische Einschr\"ankung dissipiert, was robuste Persistenz erm\"oglicht.
    \item \textbf{Endliche Negativit\"at als Eichsignal:} Genetischer Konflikt, Krebs und Wettbewerb sind das biologische Analogon der endlichen Negativit\"at im arithmetischen Kern. Sie sind keine zuf\"alligen Defekte, sondern Signale einer fehlenden organisatorischen Eichung---die durch die Entstehung neuer Kooperationsebenen (Haupt\"uberg\"ange) aufgel\"ost wird.
    \item \textbf{Eingeschr\"ankte funktionale Positivit\"at (Biologische Attraktoren):} Die stabilen Attraktoren eines Gennetzwerks oder \"Okosystems sind die ``eichstabilen'' Zust\"ande, in denen das Freie-Energie-Funktional ein eingeschr\"anktes Minimum erreicht. So wie die RH auf eine Rang-2-Stabilisierung reduziert wird, ist ph\"anotypische Stabilit\"at das Ergebnis einer endlichen Menge epigenetischer Einschr\"ankungen, die globale thermodynamische Positivit\"at gew\"ahrleisten.
\end{enumerate}

Diese Ausrichtung erhebt die Biologische Spieltheorie von einer interdisziplin\"aren Inspiration zu einer strukturellen Analogie zu ``Funktionale Positivit\"at unter Einschr\"ankung''.

\paragraph{Terminologie.} Im gesamten Paper bezeichnet \emph{resolvent-stabil} einen Fixpunkt der Dynamik, bei dem die St\"orungsanalyse zweiter Ordnung (Hessian-Eigenwerte oder Resolvent-Entwicklungskoeffizienten) in allen Nicht-Eich-Richtungen strikt positive Werte liefert. Dies ist eine strukturelle Analogie zur spektralen L\"uckeneigenschaft des FST-RH-Frameworks \citep{FST-RH2}; die Anwendung auf biologische Systeme ist konzeptuell und erfordert dom\"anenspezifische Verifikation.

% ===================================================================
\section{Thermodynamische Grundlagen: MEPP und dissipative Adaptation}
\label{sec:thermodynamics}
% ===================================================================

\subsection{Leben als Nicht-Gleichgewichts-Ph\"anomen}

Die Einordnung biologischer Prozesse in der Thermodynamik beginnt mit der Erkenntnis, dass Leben ein Ph\"anomen weit vom thermodynamischen Gleichgewicht ist \citep{EcosystemMEPP}. W\"ahrend geschlossene Systeme unvermeidlich zum Zustand maximaler Entropie und W\"armetod tendieren, erhalten sich biologische Systeme durch kontinuierlichen Export von Entropie an ihre Umgebung---ein Prozess, den Erwin Schr\"odinger 1944 als den Erwerb von ``Negentropie'' beschrieb \citep{EcosystemMEPP}.

\subsection{Drei thermodynamische Hypothesen}

Innerhalb des FST-III-Rahmens beschreiben drei zentrale thermodynamische Hypothesen das Verhalten lebender Materie \citep{EcosystemMEPP}:

\begin{enumerate}
    \item \textbf{Thermodynamischer Fingerabdruck des Lebens:} Lebende Gemeinschaften erh\"ohen die Rate der Entropieproduktion signifikant im Vergleich zu abiotischen Gegenst\"ucken unter identischen Randbedingungen.

    \item \textbf{Zustandsselektionsprinzip (Stabilit\"atsinterpretation):} Wenn ein System mehrere stabile Zust\"ande erreichen kann, tendiert die resolvent-stabile Konfiguration -- bei der Destabilisierungskan\"ale zweiter Ordnung unterdr\"uckt sind -- dazu, hohe (wenn auch nicht notwendigerweise maximale) Entropieproduktionsraten aufzuweisen.

    \item \textbf{Gradientenantwortprinzip:} Wenn der \"au\ss{}ere thermodynamische Gradient zunimmt, geht das System in einen neuen stabilen Zustand \"uber, der durch noch h\"ohere Entropieproduktion charakterisiert ist.
\end{enumerate}

\begin{table}[ht]
\centering
\caption{Thermodynamische Prinzipien und ihre biologische Relevanz. MEPP gilt f\"ur stark getriebene Systeme; Prigogines minimale Entropieproduktion gilt nur nahe dem Gleichgewicht.}
\label{tab:thermo-principles}
\begin{tabular}{lll}
\toprule
\textbf{Prinzip} & \textbf{Beschreibung} & \textbf{Biologische Relevanz} \\
\midrule
MEPP & Selektion des schnellsten Dissipationspfads & Metabolische Effizienzselektion \\
Zustandsselektion & Stabilit\"at korreliert mit hohem $\dot{S}$ & Stabilisierung komplexer Nahrungsnetze \\
Gradientenantwort & Adaptation an zunehmenden Energieeintrag & Erh\"ohung photosynthetischer Effizienz \\
Prigogine (LMEP) & Minimale Entropie nahe Gleichgewicht & Nur f\"ur lineare, schwach getriebene Systeme \\
\bottomrule
\end{tabular}
\end{table}

\begin{remark}[MEPP als Stabilit\"atsfilter, nicht als Selektionsgesetz]
\label{rem:mepp-stability-biology}
Gem\"a\ss{} der Resolventenanalyse aus dem Riemannschen-Vermutungs-Programm \citep{FST-RH2} m\"ussen die drei obigen thermodynamischen Hypothesen als \emph{Stabilit\"atsfilter} gelesen werden, nicht als Selektionsgesetze. MEPP treibt biologische Systeme nicht kausal in Richtung Zust\"ande maximaler Entropieproduktion. Stattdessen ist die Dynamik f\"uhrender Ordnung entartet: viele biologische Konfigurationen sind in erster Ordnung lebensfähig. MEPP wirkt als Filter zweiter Ordnung, der resolventeninstabile Konfigurationen eliminiert---solche, die anf\"allig f\"ur st\"orungsinduzierte Zusammenbr\"uche durch Kreuzkopplung zwischen metabolischen, strukturellen und replikativen Freiheitsgraden sind. Der Zustand mit ``h\"ochster Entropieproduktion'' wird pr\"aziser als der einzige resolventenstabile Fixpunkt innerhalb der entarteten Mannigfaltigkeit thermodynamisch lebensfähiger Konfigurationen charakterisiert. Diese Verfeinerung adressiert die Gutachter-Kritik, dass MEPP als Selektionsgesetz verwendet wird: es ist ein Stabilit\"atskriterium, kein kausaler Treiber.
\end{remark}

\subsection{Dissipative Adaptation: Die Physik der Selbstorganisation}

Ein entscheidender Br\"uckenschlag zwischen reiner Thermodynamik und biologischer Struktur wird durch Jeremy Englands Theorie der dissipativen Adaptation geliefert \citep{EnglandJCP,EnglandNano}. England argumentiert, basierend auf Nicht-Gleichgewichtsstatistik, dass Materie, die einer starken externen Energiequelle ausgesetzt ist, dazu tendiert, sich in Konfigurationen zu organisieren, die besonders gut darin sind, Energie zu absorbieren und als W\"arme zu dissipieren \citep{EnglandJCP}.

Der Mechanismus der dissipativen Adaptation beruht auf der Irreversibilit\"at von Zustands\"uberg\"angen. Wenn ein System Energie aus einer externen Quelle absorbiert, kann es energetische Aktivierungsbarrieren \"uberwinden, die durch rein thermische Fluktuationen un\"uberwindbar w\"aren. Sobald diese Energie dissipiert ist, fehlt dem System die Energie, um zur\"uck zum Anfangszustand zu springen. \"Uber die Zeit akkumulieren sich Strukturen, die hohe Resonanz mit der externen Energiequelle aufweisen \citep{EnglandNano}.

Dieser Prozess erkl\"art die spontane Entstehung selbstreplizierender Strukturen: Replikation ist eine der effizientesten Methoden, um die Kapazit\"at eines Systems zur Energieabsorption und -dissipation exponentiell zu steigern \citep{EnglandNano}. In diesem Kontext erscheint ``Fitness'' nicht als rein biologisches Konzept, sondern als physikalische Eigenschaft struktureller Persistenz unter energetischem Antrieb \citep{EnglandJCP}.

% ===================================================================
\section{Spieltheoretische Stabilit\"at: Molekulare ESS}
\label{sec:molecular}
% ===================================================================

Sobald physikalische Strukturen die Schwelle zur biologischen Reproduktion \"uberschreiten, werden ihre Interaktionen durch spieltheoretische Regeln bestimmt.

\subsection{Nash-Gleichgewicht und ESS}

Das zentrale Konzept ist das Nash-Gleichgewicht, 1950 von John Nash definiert, das Zust\"ande beschreibt, in denen kein Akteur seinen Payoff durch einseitige \"Anderung seiner Strategie verbessern kann \citep{NashPNAS,NashSFI}. In der Biologie entwickelten John Maynard Smith und George Price dies zur Evolutionarily Stable Strategy (ESS) \citep{MaynardSmithPrice1973,MaynardSmith1982}.

\begin{definition}[Evolutionarily Stable Strategy]
Eine ESS ist eine Strategie, die, wenn sie von den meisten Mitgliedern einer Population angenommen wird, nicht von einer seltenen Mutanten-Strategie invadiert werden kann. Jede ESS ist ein Nash-Gleichgewicht, aber nicht jedes Nash-Gleichgewicht ist eine ESS, da ESS zus\"atzliche Stabilit\"atsbedingungen gegen Invasionen erfordert \citep{MaynardSmith1982}.
\end{definition}

In FST-III wird dieses Prinzip auf die molekulare Ebene \"ubertragen. Proteine, Enzyme und regulatorische RNAs werden als ``Spieler'' betrachtet, deren ``Strategien'' aus ihren strukturellen Konfigurationen und katalytischen Aktivit\"aten bestehen \citep{GameBiochemistry}.

\begin{remark}[ESS als resolventenstabiler Fixpunkt]
\label{rem:ess-resolvent}
Eine Schl\"usseleinsicht aus der Resolventenanalyse der Riemannschen Vermutung \citep{FST-RH2} verfeinert die Interpretation von ESS. Eine ESS ist \emph{kein} globales Fitnessoptimum; sie ist ein resolventenstabiler Fixpunkt. In f\"uhrender Ordnung ist die Fitnesslandschaft entartet---viele ph\"anotypische Konfigurationen ergeben vergleichbare Payoffs. Die Selektion erfolgt in zweiter Ordnung durch resolventenartige Kopplungen zwischen den entarteten Modi: Invasionsdynamik, frequenzabh\"angige Interaktionen und r\"aumliche Kopplung spalten die Entartung gemeinsam auf. Die ESS ist die einzige Konfiguration, bei der alle Destabilisierungskan\"ale zweiter Ordnung gleichzeitig unterdr\"uckt sind. Nash-Frustration (Abschnitt~\ref{sec:molecular}) ist kein Defekt, sondern ein \emph{notwendiges strukturelles Merkmal} dieser Resolventenarchitektur: sie reflektiert die residuale Kopplung zweiter Ordnung, die die Konfiguration stabilisiert und gleichzeitig eine Kristallisation in ein rigides, fragiles Optimum verhindert.
\end{remark}

\begin{table}[ht]
\centering
\caption{Abbildung spieltheoretischer Konzepte auf physikalische und biologische Gegenst\"ucke.}
\label{tab:game-biology-mapping}
\begin{tabular}{lll}
\toprule
\textbf{Spielkonzept} & \textbf{Physikalische Entsprechung} & \textbf{Biologisches Beispiel} \\
\midrule
Nash-Gleichgewicht & Stabiler station\"arer Zustand & Ph\"anotyp-Koexistenz \\
ESS & Attraktor mit lokaler Stabilit\"at & Allel-Fixierung \\
Payoff-Matrix & Energie-/Fitnesslandschaft & Enzym-Substrat-Affinit\"at \\
Gemischte Strategie & Polymorphismus/Heterogenit\"at & Arbeitsteilung in Kolonien \\
\bottomrule
\end{tabular}
\end{table}

\subsection{Proteinfaltung als Optimierungsspiel}

Ein markantes Beispiel ist die Proteinfaltung, die als Spiel gegen die thermodynamische Landschaft modelliert werden kann \citep{EnzymeFolding}. Die Aminos\"auresequenz versucht, ein globales Minimum der freien Energie zu erreichen, was einem Nash-Gleichgewicht in einem n-Personen-Spiel molekularer Kr\"afte entspricht:

\begin{itemize}
    \item \textbf{Spieler:} Einzelne Aminos\"aurereste
    \item \textbf{Strategien:} R\"uckgrat-Diederwinkel $(\phi, \psi)$
    \item \textbf{Payoff:} $-\Delta G$ (negative Freie-Energie-\"Anderung)
    \item \textbf{Nash-Gleichgewicht:} Native Konformation (Anfinsens Dogma \citep{Anfinsen1973})
\end{itemize}

Fehlgefaltete Proteine (Prionen) repr\"asentieren alternative Gleichgewichte---Konfigurationen, bei denen keine einseitige Abweichung den Payoff verbessert, obwohl der globale Payoff suboptimal ist.

\subsection{Nash-Frustration vs.\ energetische Frustration}

Das Konzept der lokalen Frustration in Proteinstrukturen wurde extensiv durch den Protein Frustratometer \citep{Ferreiro2007,Ferreiro2014} untersucht, der Reste identifiziert, bei denen die native Aminos\"aureidentit\"at suboptimal ist, gegeben die lokale strukturelle Umgebung. Diese energetische Frustration korreliert mit funktionellen Stellen, Bindungsschnittstellen und allosterischen Regionen.

Der FST-III-Rahmen f\"uhrt ein komplement\"ares Konzept ein: \textbf{Nash-Frustration}. W\"ahrend energetische Frustration fragt ``welche Reste haben suboptimale lokale Energie?'', fragt Nash-Frustration ``welche Reste verhindern spieltheoretische Koordination?''

Formal wird Nash-Frustration aus dem Best-Response-Jacobian $J = I - \eta H$ berechnet, wobei $H$ die Hesse-Matrix des Potentials und $\eta > 0$ ein Lernraten-Parameter ist. Der Frustrationswert f\"ur Rest $i$ ist:
\begin{equation}
\text{frustration}(i) = \sum_{k: |\lambda_k| \geq 1} (|\lambda_k| - 1) \left( v_{k,\phi_i}^2 + v_{k,\psi_i}^2 \right)
\end{equation}
wobei $\lambda_k$ und $v_k$ Eigenwerte und Eigenvektoren von $J$ sind.

\textbf{Wichtige Einschr\"ankung:} Die Lernrate $\eta$ ist ein freier Parameter, der direkt den Spektralradius $\rho(J)$ und damit das Ausma\ss{} der Nash-Frustration bestimmt. Der Wert $\eta = 0{,}01$, der in der HP35-Berechnung unten verwendet wird, wurde gew\"ahlt, um die native Faltung nahe der Stabilit\"atsgrenze $\rho(J) \approx 1$ zu platzieren; andere Wahlen von $\eta$ ergeben qualitativ unterschiedliche Frustrationskarten. Eine prinzipielle, datengetriebene Methode zur Auswahl von $\eta$---beispielsweise die Kalibrierung gegen bekannte Faltungskinetiken oder allosterische \"Ubergangsraten---ist eine notwendige Voraussetzung, bevor Nash-Frustrationswerte \"uber Proteine hinweg verglichen oder als Pr\"adiktoren verwendet werden k\"onnen. Die hier berichteten Ergebnisse sollten daher als Machbarkeitsstudie interpretiert werden, nicht als validiertes quantitatives Werkzeug.

\begin{table}[ht]
\centering
\caption{Vergleich von energetischer Frustration (Ferreiro et al.) und Nash-Frustration (diese Arbeit). Die beiden Ma\ss{}e erfassen komplement\"are Aspekte der Proteinstabilit\"at.}
\label{tab:frustration-comparison}
\begin{tabular}{lll}
\toprule
\textbf{Konzept} & \textbf{Energetische Frustration} & \textbf{Nash-Frustration} \\
\midrule
Frage & Ist lokale Energie optimal? & Ist einseitige Abweichung stabil? \\
Methode & Kontaktenergie-Statistik & Best-Response-Jacobian-Eigenanalyse \\
Interpretation & Evolution\"are Einschr\"ankung & Spieltheoretische Koordination \\
Hohe Werte zeigen & Funktionelle Stellen, Schnittstellen & Faltungsintermediate, kinetische Fallen \\
\bottomrule
\end{tabular}
\end{table}

Rechnerische Experimente an drei strukturell verschiedenen Proteinen best\"atigen, dass Nash-Frustration biologisch bedeutsame Reste identifiziert. Tabelle~\ref{tab:nash-results} fasst die Ergebnisse zusammen.

\begin{table}[ht]
\centering
\caption{Nash-Gleichgewichtsanalyse von Proteinfaltungen. Parameter gelernt aus HP35 (Villin Headpiece, 1YRF) und angewendet zur Vorhersage der Faltung von Protein G (1PGA, 56 Reste) und der p53-DNA-Bindungsdom\"ane (2XWR, 462 Reste). Best-Response-Dynamik konvergiert \"uber alle drei Proteine trotz einer Gr\"o\ss{}enordnung Unterschied in der Proteingr\"o\ss{}e. Die Trainingsproteindaten (Spektralradius, Frustrationslandkarte) charakterisieren die Nash-Stabilit\"atslandschaft.}
\label{tab:nash-results}
\begin{tabular}{lccccc}
\toprule
\textbf{Protein} & \textbf{Reste} & $\phi_{\text{rms}}$ [rad] & $\psi_{\text{rms}}$ [rad] & \textbf{Konvergiert} & $F_{\text{best}}$ \\
\midrule
HP35 (1YRF, Training) & 35 & 1{,}054 & 1{,}661 & Ja (267 Sw) & $-4{,}61$ \\
Protein G (1PGA) & 56 & 1{,}300 & 2{,}002 & Ja (282 Sw) & $-10{,}61$ \\
p53 DBD (2XWR) & 462 & 1{,}247 & 2{,}114 & Ja & $-39{,}43$ \\
\bottomrule
\end{tabular}
\end{table}

\textbf{Trainingsprotein (HP35).} Nash-Frustration konzentriert sich an Loop-Regionen (insbesondere Met11, Pro20 mit Scores 1{,}0 und 0{,}49), w\"ahrend Helix-Kerne (Ser14, Asn18, Trp22) nahezu Null-Frustration aufweisen---konsistent mit ihrer Rolle als ``strukturelle Anker'' im Faltungsspiel. Der Spektralradius $\rho(J) = 1{,}013$ mit 80\% positiven Hessian-Eigenwerten und mittlerer Frustration 0{,}082 charakterisiert eine Nahe-Nash-Struktur. Dieses Muster unterscheidet sich von energetischer Frustration, die typischerweise an oberfl\"achenexponierten funktionellen Resten ihren H\"ochstwert erreicht.

\textbf{Cross-Protein-Vorhersage.} Die aus HP35 gelernten Modellparameter wurden verwendet, um die R\"uckgrat-Winkel von Protein G (1PGA, 56 Reste) und der p53-Tumorsuppressor-DNA-Bindungsdom\"ane (2XWR, 462 Reste) mittels Best-Response-Dynamik vorherzusagen. Beide Vorhersagen konvergierten zu wohldefinierten Strukturen (Tabelle~\ref{tab:nash-results}). Bemerkenswert ist, dass die p53-Dom\"ane---13$\times$ gr\"o\ss{}er als das Trainingsprotein und einer v\"ollig anderen Faltungsklasse zugeh\"orig (Immunglobulin-artiges $\beta$-Sandwich vs.\ $\alpha$-helikal)---einen $\phi_{\text{rms}}$ von 1{,}25\,rad erreichte, vergleichbar mit Protein G (1{,}30\,rad). W\"ahrend die absoluten RMSD-Werte gro\ss{} sind (erwartbar f\"ur ein 70-Parameter-Modell ohne Seitenketten-Terme), ist die Schl\"usselbeobachtung, dass Best-Response-Dynamik \emph{\"uber alle drei Proteine konvergiert}---was darauf hindeutet, dass die Potentialspiel-Formulierung universelle Merkmale der kooperativen Logik der Proteinfaltung erfasst.

\paragraph{Rechnerische Beobachtung (Nash-Stabilit\"atsgrenze).}
Native Proteinfaltungen sind keine exakten Nash-Gleichgewichte unter lokaler Best-Response-Dynamik, sondern liegen nahe der Stabilit\"atsgrenze ($\rho(J) \approx 1$). Der Spektralradius $\rho(J) = 1{,}013$ f\"ur HP35 zeigt, dass die native Struktur nahezu selbststabilisierend ist, mit 14 von 35 instabilen Moden, die nur an funktionellen Scharnierregionen lokalisiert sind, statt globalem Koordinationsversagen.

Dieses Ergebnis hat tiefgreifende Implikationen: Nash-Stabilit\"at ist \textit{strikter} als Energieminimierung. Eine Struktur kann ein lokales Energieminimum sein, w\"ahrend sie dennoch einseitige Verbesserungen f\"ur einzelne Reste erlaubt. Die lokalisierte Nash-Frustration bei Met11 identifiziert einen funktionellen Freiheitsgrad---ein ``Scharnier'', das konformationelle Flexibilit\"at erm\"oglicht, w\"ahrend die Helix-Kerne stabile Koalitionen bilden. Dies ist pr\"azise das, was man f\"ur funktionelle Proteine erwartet: sie balancieren Stabilit\"at gegen Flexibilit\"at, anstatt in rigide Nash-Optima zu kristallisieren. Eine systematische Kalibrierung von $\eta$ gegen experimentelle Daten---etwa durch Abgleich der Nash-instabilen Reste mit bekannten allosterischen Stellen oder Faltungskeimen aus H/D-Austausch-Experimenten---w\"urde diesen Proof-of-Concept in ein pr\"adiktives Werkzeug transformieren. Bis eine solche Kalibrierung durchgef\"uhrt ist, sollten die hier berichteten numerischen Werte ($\rho(J) = 1{,}013$, Frustrations-Maxima bei Met11 und Pro20) als parameterabh\"angige Illustrationen statt als robuste quantitative Vorhersagen interpretiert werden.

\subsection{Genregulation als Signalspiel}

Genregulationsnetzwerke (GRNs) steuern Genexpressionsmuster und k\"onnen als Signalspiele modelliert werden \citep{PredatorPreyESS}:

\begin{itemize}
    \item \textbf{Spieler:} Transkriptionsfaktoren und DNA-Bindungsstellen
    \item \textbf{Strategien:} Binden/Nicht-Binden
    \item \textbf{Payoff:} Organismische Fitness
\end{itemize}

Die mathematische Verbindung zwischen Spieltheorie und Populationsdynamik wird durch die Replikatorgleichung hergestellt, die formal \"aquivalent zu Lotka-Volterra-Modellen ist \citep{CoexistenceLotkaVolterra}. Ein stabiler Fixpunkt in diesen Gleichungen entspricht einer ESS und somit einem Verhaltensattraktor, der die Langzeitstruktur des Systems bestimmt.

% ===================================================================
\section{Das Free Energy Principle und Active Inference}
\label{sec:fep}
% ===================================================================

Das Free Energy Principle (FEP), prim\"ar entwickelt von Karl Friston, liefert eine vereinigende Theorie f\"ur biologische Selbstorganisation, die Informationstheorie und Thermodynamik verschmilzt \citep{FristonFreeEnergy,FristonPDF}. \textbf{Hinweis:} Die Behauptung, dass FEP und MEPP kompatibel oder \"aquivalent sind, ist nicht unumstritten. FEP minimiert variationelle freie Energie (eine informationstheoretische Gr\"o\ss{}e, die Surprise begrenzt), w\"ahrend MEPP die thermodynamische Entropieproduktionsrate maximiert. Die Verbindung zwischen diesen beiden Prinzipien ist ein aktives Diskussionsgebiet \citep{FristonFreeEnergy}; das vorliegende Paper behandelt sie als komplement\"are Perspektiven auf verschiedenen Analyseebenen, leitet aber nicht formal eines aus dem anderen ab.

\subsection{Theoretischer Rahmen}

FEP postuliert, dass jedes biologische System, das sich erfolgreich gegen Entropie (Zerfall) verteidigt, seine ``variationelle freie Energie'' minimieren muss \citep{FristonFreeEnergy}. Diese Gr\"o\ss{}e repr\"asentiert eine mathematische obere Schranke der ``Surprise'' (Surprisal), die das System durch sensorische Inputs erlebt \citep{FristonFreeEnergy}.

Ein biologischer Agent existiert innerhalb einer ``Markov-Decke''---einer statistischen Grenze, die interne Zust\"ande von externen Zust\"anden trennt \citep{FristonPDF}. Um zu \"uberleben, muss der Agent sicherstellen, dass seine sensorischen Zust\"ande innerhalb physiologischer Grenzen bleiben. Dies wird durch zwei Mechanismen erreicht:

\begin{enumerate}
    \item \textbf{Wahrnehmung:} Das System passt seine internen Modelle an, um die Ursachen sensorischer Daten besser vorherzusagen (Minimierung des Vorhersagefehlers) \citep{FristonFreeEnergy}
    \item \textbf{Active Inference:} Das System handelt, um die Welt zu ver\"andern oder sensorische Daten zu selektieren, die seinen internen Erwartungen entsprechen \citep{FristonFreeEnergy}
\end{enumerate}

Durch Minimierung der freien Energie wird das System implizit zu einem Modell seiner Umgebung. Dieser Prozess ist eng mit MEPP verkn\"upft: w\"ahrend das System intern seine freie Energie minimiert (Ordnung schafft), maximiert es die Entropieproduktion in der Umgebung durch metabolische und agentive Aktivit\"aten \citep{FristonPDF}.

In FST-III sind Verhaltensattraktoren diejenigen Zust\"ande, in denen die freie Energie minimal ist und der Agent eine stabile, vorhersagbare Interaktion mit seiner Nische erreicht hat \citep{FristonFreeEnergy}.

\subsection{Freie-Energie-Minimierung als Nash-Fixpunkt}

Die Verbindung zwischen FEP und Nash-Gleichgewicht kann unter bestimmten Bedingungen formalisiert werden. Sei $p(o,s_1,\dots,s_N)$ ein gemeinsames generatives Modell f\"ur $N$ interagierende Agenten, wobei jeder Agent $i$ eine faktorisierte approximative Posterior $q_i(s_i)$ \"uber seine latenten Zust\"ande kontrolliert. Die gemeinsame variationelle freie Energie ist:
\begin{equation}
F(q; o) = \mathbb{E}_q\!\left[\ln q(s) - \ln p(o, s)\right], \qquad q(s) = \prod_{i=1}^{N} q_i(s_i).
\end{equation}

Definiere ein Spiel, in dem jeder Agent $i$ die Kontrolle \"uber $q_i$ hat und alle Agenten das Kostenfunktional $J_i(q_i, q_{-i}) := F(\prod_j q_j; o)$ teilen. Dann ist die beste Antwort von Agent $i$ gegen festes $q_{-i}$ pr\"azise das koordinatenweise variationelle Update (Mean-Field-VI), und ein Nash-Gleichgewicht $q^* = (q_1^*, \dots, q_N^*)$ ist ein Fixpunkt verteilter Freie-Energie-Minimierung. Dies konstituiert ein \emph{Potentialspiel} mit Potential $F$: Nash-Gleichgewichte fallen mit koordinatenweisen Minima der gemeinsamen freien Energie zusammen.

\paragraph{G\"ultigkeitsbereich der Br\"ucke.} Diese formale \"Aquivalenz gilt f\"ur die Belief-Optimierungskomponente von FEP (variationelle Inferenz). F\"ur Active Inference, wo Agenten zus\"atzlich Policies $\pi_i$ selektieren, um die erwartete freie Energie $G(\pi_i, \pi_{-i})$ zu minimieren, wird die Nash-Bedingung: kein Agent kann einseitig seine Policy \"andern, um seine erwartete freie Energie zu reduzieren---ein genuines Mehr-Agenten-Gleichgewicht, das Interaktion durch Zustandsdynamik involviert. Das vorliegende Paper entwickelt diese Erweiterung nicht formal; der Beitrag besteht darin, die pr\"azise Demarkation zu liefern: FEP-als-Nash ist mehr als Metapher, wenn das Spiel explizit als variationelles/Potentialspiel \"uber Verteilungen definiert ist, und blo\ss{} konzeptuell, wenn ``Nash'' als Synonym f\"ur ``Fixpunkt'' verwendet wird.

\subsection{Skalentrennung: Aufl\"osung des FEP-MEPP-Paradoxons}
\label{sec:fep-mepp-paradox}

Eine persistente konzeptuelle Spannung in der theoretischen Biologie fragt, wie ein System seine interne freie Energie \emph{minimieren} (FEP) kann, w\"ahrend es gleichzeitig die Entropieproduktion in der Umgebung \emph{maximiert} (MEPP). Die folgende Proposition l\"ost dieses Paradoxon durch Skalentrennung \"uber die Markov-Decke.

\begin{proposition}[FEP-MEPP-Dualit\"at via Markov-Decke (Konjektur)]
\label{prop:fep-mepp}
Sei $\mathcal{B}$ ein selbstorganisierendes System mit Markov-Decke $\partial\mathcal{B}$, internen Zust\"anden $s$ und Umgebungszust\"anden $e$. Definiere:
\begin{itemize}
  \item $F_{\mathrm{int}}(q) :=
    \mathbb{E}_{q}\!\bigl[\ln q(s) - \ln p(o, s)\bigr]$,
    die variationelle freie Energie des internen Modells;
  \item $\dot{S}_{\mathrm{ext}} :=
    \frac{\dd S_{\mathrm{env}}}{\dd t} =
    \frac{1}{T}\,\frac{\dd Q}{\dd t}$,
    die Rate der Entropieproduktion in der Umgebung.
\end{itemize}
Wenn $\mathcal{B}$ $F_{\mathrm{int}}$ minimiert (FEP), indem es ein akkurates internes Modell aufrechterh\"alt und Aktionen w\"ahlt, die Surprise reduzieren, dann:
\begin{enumerate}
  \item[\emph{(i)}] Die strukturelle Integrit\"at von $\mathcal{B}$ wird aufrechterhalten: $\mathcal{B}$ widersteht der Aufl\"osung und erh\"alt die Markov-Decke (niedrige interne Entropie);
  \item[\emph{(ii)}] Die intakte, fernab-vom-Gleichgewicht-Struktur $\mathcal{B}$ fungiert als effiziente dissipative Maschine, die Energiegradienten aus der Umgebung durch metabolische und agentive Aktivit\"aten kanalisiert;
  \item[\emph{(iii)}] Unter allen Randbedingungen, die konsistent mit dem \"Uberleben von $\mathcal{B}$ sind, maximiert die tats\"achliche Dynamik die Rate der Entropieproduktion $\dot{S}_{\mathrm{ext}}$ in der Umgebung (MEPP-Selektion).
\end{enumerate}
Kurz: \emph{interne Freie-Energie-Minimierung erm\"oglicht externe Entropieproduktions-Maximierung}. Die beiden Prinzipien stehen nicht im Widerspruch, sondern operieren auf komplement\"aren Seiten der Markov-Decke.
\end{proposition}

\noindent\textbf{Status:} Diese Proposition ist als Konjektur formuliert, gest\"utzt auf qualitative Argumente. Ein formaler Beweis w\"urde erfordern, dass FEP-minimierende Systeme notwendigerweise die Entropieproduktion in der Umgebung unter den genannten Randbedingungen maximieren -- ein Resultat, das in der Literatur nicht demonstriert wurde.

\begin{remark}[Spieltheoretische Interpretation]
In der Sprache des vorliegenden Papers entspricht die FEP-MEPP-Dualit\"at einem Zweistufen-Spiel:
\begin{itemize}
  \item \emph{Inneres Spiel} (innerhalb $\partial\mathcal{B}$): Agenten minimieren die variationelle freie Energie $F_{\mathrm{int}}$ \"uber den Nash-Fixpunkt von Proposition~\ref{prop:fep-mepp}---das oben beschriebene Potentialspiel.
  \item \emph{\"Au\ss{}eres Spiel} (\"uber $\partial\mathcal{B}$): Das Ensemble \"uberlebender Organismen konkurriert um Energiegradienten. Organismen, die effizienter dissipieren (h\"oheres $\dot{S}_{\mathrm{ext}}$), verdrängen langsamere Dissipatoren, wodurch MEPP als makroskaliges Selektionskriterium etabliert wird.
\end{itemize}
Dies l\"ost das von Martyushev~\cite{Martyushev2010} aufgeworfene MEPP-Paradoxon: MEPP gilt f\"ur das \emph{umgebungsgekoppelte} System, nicht f\"ur die internen Freiheitsgrade, wo minimale Entropieproduktion (Prigogine) lokal gelten kann. Die Nash-Frustration $\rho(J)$ aus Abschnitt~\ref{sec:molecular} misst die residuale Spannung zwischen innerer Optimalit\"at und \"au\ss{}erem dissipativem Druck.
\end{remark}

% ===================================================================
\section{Boolesche Netzwerke und der Rand des Chaos}
\label{sec:boolean}
% ===================================================================

Auf der Ebene der Genregulation manifestiert sich Stabilit\"at in der Dynamik Boolescher Netzwerke, wie sie von Stuart Kauffman eingef\"uhrt wurden \citep{BooleanPost,KauffmanOrigins}.

\subsection{Netzwerkdynamik}

Diese Netzwerke modellieren Gene als bin\"are Schalter, deren Zust\"ande durch regulatorische Logik verkn\"upft sind \citep{BooleanPost}. Die Langzeitdynamik solcher Netzwerke flie\ss{}t in Attraktoren, die biologisch als spezifische Zelltypen oder funktionelle Zust\"ande interpretiert werden \citep{BooleanEvolution}.

Kauffman entdeckte, dass Netzwerkstabilit\"at von der Konnektivit\"at und dem Typ der Booleschen Funktionen abh\"angt \citep{BooleanPost}. Netzwerke am ``Rand des Chaos''---einer kritischen \"Ubergangsregion zwischen rigider Ordnung und unvorhersagbarem Chaos---weisen die h\"ochste Evolvabilit\"at und Informationsverarbeitungskapazit\"at auf:

\begin{table}[ht]
\centering
\caption{Regime Boolescher Netzwerke und ihre biologische Interpretation. Netzwerke am Rand des Chaos weisen maximale Evolvabilit\"at auf.}
\label{tab:network-regimes}
\begin{tabular}{lll}
\toprule
\textbf{Netzwerkregime} & \textbf{Dynamisches Verhalten} & \textbf{Biologische Funktion} \\
\midrule
Geordnet & Eingefrorene Zust\"ande, geringe Variabilit\"at & Strukturelle Integrit\"at, Basalmetabolismus \\
Kritisch (Rand des Chaos) & Komplexe Muster, hoher Informationsfluss & Adaptives Verhalten, Differenzierung, Lernen \\
Chaotisch & Hohe Sensitivit\"at, Instabilit\"at & Pathologische Zust\"ande, Malignit\"at \\
\bottomrule
\end{tabular}
\end{table}

Diese Attraktoren in Genregulationsnetzwerken bilden die materielle Basis f\"ur die spieltheoretischen Strategien auf h\"oheren Ebenen. Eine ESS auf der Makroebene erfordert stabile genetische Attraktoren auf der Mikroebene, die zuverl\"assig die notwendigen ph\"anotypischen Merkmale produzieren \citep{BooleanMicroarray,KauffmanOrigins}.

% ===================================================================
\section{Zellul\"are Ebene: Multizellularit\"at als Koalitionsspiel}
\label{sec:multicellularity}
% ===================================================================

\subsection{Das Kooperationsproblem}

In einem multizellul\"aren Organismus ``opfern'' einzelne Zellen ihr Reproduktionspotential zugunsten des Ganzen. Somatische Zellen reproduzieren sich nicht direkt; nur Keimbahnzellen \"ubertragen Gene.

Dies repr\"asentiert extreme Kooperation: Zellen adoptieren eine Strategie (Differenzierung, Nicht-Reproduktion), die dem Kollektiv auf individuelle Kosten n\"utzt.

\subsection{Koalitionsspiel-Formulierung}

\begin{itemize}
    \item \textbf{Spieler:} Einzelne Zellen innerhalb des Organismus
    \item \textbf{Strategien:} Kooperieren (Entwicklungsprogramm folgen) oder Defektieren (autonom proliferieren)
    \item \textbf{Payoff:} Genetische Fitness (\"uber Keimbahn \"ubertragen)
\end{itemize}

Alle Zellen in einem Organismus sind genetisch identisch (klonal). Daher ist der Payoff f\"ur jede Zelle gleich dem Payoff der Keimbahn: wenn sich der Organismus reproduziert, werden die Gene aller Zellen \"ubertragen.

\begin{proposition}[Multizellularit\"at als Nash-Gleichgewicht]
\label{prop:multicellularity}
In einem klonalen Organismus ist Kooperation ein Nash-Gleichgewicht, weil Defektion die Fitness genetisch identischer Zellen reduziert. Dies ist eine Reformulierung von Hamiltons Regel \citep{Hamilton1964} mit $r=1$ in spieltheoretischer Sprache; der Beitrag hier besteht in der Einbettung in den hierarchischen FST-Rahmen.
\end{proposition}

\subsection{Krebs als Koalitionszusammenbruch}

Krebs als Zusammenbruch multizellul\"arer Kooperation ist ein gut etablierter Rahmen \citep{AktipisCancer}. Die spieltheoretische Rahmung hier folgt Aktipis et al.\ (2015), die evolution\"are Spieltheorie anwenden, um somatische Evolution und Krebsprogression zu erkl\"aren. Der FST-III-Beitrag besteht darin, diesen Rahmen in das hierarchische Nash-Gleichgewichts-Bild einzubetten: Krebs repr\"asentiert nicht blo\ss{} einen selektiven Vorteil f\"ur einzelne Zellen, sondern einen Phasen\"ubergang auf der Koalitionsebene, analog zum Parasitenproblem in der pr\"abiotischen Chemie, das in FST-II diskutiert wird. Formal repr\"asentiert Krebs einen Koalitionszusammenbruch:
\begin{enumerate}
    \item Somatische Mutationen erzeugen genetische Divergenz
    \item Mutierte Zellen teilen nicht mehr alle Gene mit anderen Zellen
    \item Defektion (Proliferation) n\"utzt nun der mutierten Linie
    \item Das Nash-Gleichgewicht bricht zusammen; Krebszellen ``defektieren''
\end{enumerate}

Dies verbindet sich mit dem Parasitenproblem in Hyperzyklen \citep{FSTII}: sowohl Krebs als auch chemische Parasiten nutzen kooperative Systeme aus, indem sie Bedingungen brechen, die Kooperation stabilisieren.

% ===================================================================
\section{Skalierung und die Renormierungsgruppe}
\label{sec:renormalization}
% ===================================================================

Ein zentrales Element von FST-III ist die Erkenntnis, dass biologische Prinzipien skaleninvariant \"uber viele Ebenen operieren \citep{ScaleRGFrontiers,ScaleRGResearch}.

\subsection{Der RG-Rahmen}

\textbf{Terminologischer Vorbehalt:} Der Begriff ``Renormierungsgruppe'' wird in diesem Abschnitt in einem erweiterten, metaphorischen Sinne verwendet. In der Physik der kondensierten Materie bezieht sich die Renormierungsgruppe (RG) auf einen pr\"azisen mathematischen Rahmen mit Flussgleichungen, Fixpunkten und Skalenexponenten \citep{ScaleRGResearch}. Die vorliegende Verwendung---die beschreibt, wie effektive spieltheoretische Parameter sich \"andern, wenn man auf sukzessive gr\"o\ss{}eren biologischen Skalen betrachtet---ruft die \emph{konzeptuelle Logik} der Grobk\"ornung auf, liefert aber keine RG-Flussgleichungen, Fixpunktanalyse oder dimensionelle Skalengesetze f\"ur biologische Systeme. Ob eine rigorose RG-Behandlung evolution\"arer \"Uberg\"ange erreichbar ist, bleibt eine offene Frage \citep{ScaleRGFrontiers}.

Mit diesem Vorbehalt im Hinterkopf: Die Grobk\"ornungsperspektive postuliert, dass biologische Systeme oft hierarchisch organisiert sind---von fraktaler Verzweigung in Blutgef\"a\ss{}en bis zu hierarchischen neuronalen Architekturen \citep{SelfSimilarNetworks}.

Renormalizing Generative Models (RGM) zeigen, dass Prinzipien der freien Energie und spieltheoretischen Optimierung rekursiv auf jeder Ebene angewandt werden \citep{ScaleRGFrontiers}. Ein Organismus kann als Grobk\"ornung \"uber Milliarden von Zellen betrachtet werden, wobei Zellen als ``Atome'' in einem Spiel h\"oherer Ordnung fungieren, dessen Regeln durch die evolution\"are Geschichte der Spezies festgelegt sind \citep{ReplicatorOptimization}.

\subsection{Das skalenrelative Kernmodell}

Diese Instanziierung wird durch das Skalenrelative Kernmodell beschrieben \citep{ReplicatorOptimization}:

\begin{table}[ht]
\centering
\caption{Skalenrelatives Kernmodell: Fitnessdefinitionen und Strategietypen auf jeder biologischen Organisationsebene.}
\label{tab:scale-kernel}
\begin{tabular}{ll}
\toprule
\textbf{Ebene} & \textbf{Fitness/Strategie} \\
\midrule
Zellul\"ar & Replikationsrate als Fitness, Mutationen als Strategiewechsel \\
Organismisch & Darwinsche Fitness, Verhaltensstrategien \\
Agentiv & Intentionalit\"at, Lernen durch Imitation \\
Institutionell & Legitimit\"at und soziale Kooperation als ESS \\
\bottomrule
\end{tabular}
\end{table}

Diese hierarchische Struktur stellt sicher, dass Konflikte zwischen Skalen (z.\,B.\ Krebs als Defektion einzelner Zellen gegen den Organismus) durch \"ubergeordnete Selektionsdr\"ucke minimiert werden \citep{ReplicatorOptimization}.

\subsection{Haupt\"uberg\"ange als RG-Schritte}

Die Haupt\"uberg\"ange der Evolution, identifiziert von Szathmary und Maynard Smith \citep{SzathmaryMaynardSmith1995}, bilden auf diesen Rahmen ab:

\begin{itemize}
    \item Gene $\to$ Chromosomen (Block 1)
    \item Prokaryoten $\to$ Eukaryoten (Block 2)
    \item Einzeller $\to$ Vielzeller (Block 3)
    \item Individuen $\to$ Gesellschaften (Block 4)
\end{itemize}

Jeder \"Ubergang schafft eine neue Ebene von ``Spielern'' aus Koalitionen niedrigerer Spieler und erh\"alt spieltheoretische Struktur durch Grobk\"ornung.

\paragraph{Ein Toy-Renormierungsschritt f\"ur r\"aumliche evolution\"are Spiele.}
Betrachte ein 2D-Gitter-Spiel mit Strategien $\sigma_x\in\{C,D\}$ und Payoff-Matrix
$A=\begin{pmatrix}R&S\\ T&P\end{pmatrix}$.
F\"ur einen $2\times 2$-Block $B$ definiere eine Blockstrategie $\Sigma_B=\mathcal B(\sigma|_B)$ durch Mehrheitsregel.
F\"ur benachbarte Bl\"ocke $B,B'$ definiere den effektiven Block-Payoff durch Randmittelung

\[
A'_{s,t}
:=\frac{1}{|\partial(B,B')|}\,
\mathbb E\!\left[\sum_{\langle x,y\rangle\in\partial(B,B')} A_{\sigma_x,\sigma_y}\ \Big|\ \Sigma_B=s,\ \Sigma_{B'}=t\right],
\qquad s,t\in\{C,D\}.
\]

In einer Frozen-Interior-Approximation, sei
$q_s=\mathbb P(\sigma_x=C\mid \Sigma_B=s,\ x\in\partial B)$.
Dann erf\"ullt die RG-Abbildung $R_2:(A,\beta)\mapsto(A',\beta')$ die explizite Mischungsformel

\[
A'_{s,t}
= q_s q_t\,R + q_s(1-q_t)\,S + (1-q_s)q_t\,T + (1-q_s)(1-q_t)\,P,
\]

mit $q_C\approx 1-\varepsilon$ und $q_D\approx \varepsilon$, kontrolliert durch die Within-Block-Dynamik (Selektionsst\"arke $\beta$ und interne Durchsetzungsmechanismen).
Fixpunkte umfassen die geordneten Phasen (All-$C$, All-$D$), w\"ahrend \"Anderungen in Within-Block-Einschr\"ankungen den RG-Fluss durch Modifikation von $q_C-q_D$ verschieben.
Dies liefert einen konkreten Sinn, in dem evolution\"are Haupt\"uberg\"ange Grobk\"ornungsschritten zusammen mit der Aktivierung neuer relevanter Operatoren (z.\,B.\ Policing, Bottlenecks) entsprechen, die kooperative Makrovariablen stabilisieren.

% ===================================================================
\section{\"Okosystem-Ebene: Das \"okologische Spiel}
\label{sec:ecosystem}
% ===================================================================

\subsection{Arten als Spieler}

Im \"okologischen Spiel:
\begin{itemize}
    \item \textbf{Spieler:} Artenpopulationen
    \item \textbf{Strategien:} \"Okologische Rollen (R\"auber, Beute, Konkurrent, Mutualist)
    \item \textbf{Payoff:} Populationswachstumsrate
\end{itemize}

\subsection{\"Okologische Nischen als Nash-Gleichgewichte}

\"Okologische Nischen sind Nash-Gleichgewichts-Konfigurationen: keine Art kann einseitig ihre Strategie \"andern, um die Fitness zu verbessern, w\"ahrend andere konstant bleiben.

\subsection{Massenaussterben als Phasen\"uberg\"ange}

Massenaussterbe-Ereignisse repr\"asentieren Phasen\"uberg\"ange im \"okologischen Spiel:
\begin{enumerate}
    \item Umweltst\"orung ver\"andert die Payoff-Matrix
    \item Vorherige Nash-Gleichgewichte werden instabil
    \item System geht in neues Gleichgewicht mit anderer Artenzusammensetzung \"uber
\end{enumerate}

Dies spiegelt kosmologische Phasen\"uberg\"ange im CFM (Cosmological Fractal Model) wider: ver\"anderte Randbedingungen destabilisieren bestehende Gleichgewichte und treiben \"Uberg\"ange zu neuen Konfigurationen.

\section{Kosmologische Erweiterungen}
\label{sec:cosmology}

Das in diesem Paper entwickelte biologisch-spieltheoretische Framework verbindet sich mit einem breiteren Forschungsprogramm, das skaleninvariante Optimierungsprinzipien \"uber Physik und Kosmologie hinweg untersucht. Spekulative Erweiterungen -- einschlie\ss{}lich Verbindungen zu modifizierter Gravitation (MOND) und kosmologischen Randbedingungen -- werden separat im FST-Framework-Paper \citep{FST-Unified2026} entwickelt und sollten nicht als etablierte Bestandteile der hier pr\"asentierten biologischen Theorie gelesen werden.

% ===================================================================
\section{Synthese: Von der Entropie zum Bewusstsein}
\label{sec:synthesis}
% ===================================================================

FST-III liefert ein integriertes Bild der Realit\"at, das sich in drei Phasen entfaltet:

\subsection{Phase 1: Thermodynamische Dissipation}

Der Energiefluss des jungen Universums zwingt einfache Materie in dissipative Strukturen. Hier herrscht MEPP; es ist die \"Ara schneller Entropieproduktion und der Entstehung von Englands dissipativen Adaptoren \citep{EcosystemMEPP}.

\subsection{Phase 2: Strategische Konsolidierung}

Diese Strukturen beginnen sich selbst zu replizieren. Wettbewerb um Ressourcen f\"uhrt zur Entstehung von Evolutionarily Stable Strategies (ESS). Die ``Spieler'' sind noch rein chemisch, aber ihre Interaktionen sind bereits spieltheoretisch beschreibbar \citep{MaynardSmithPrice1973}.

\subsection{Phase 3: Kognitive Inferenz}

Organismen entwickeln interne Modelle ihrer Umgebung, um ihre freie Energie zu minimieren. Die Entstehung der Markov-Decke erlaubt die Trennung von Selbst und Welt und f\"uhrt zur Bildung von Verhaltensattraktoren, die wir Instinkte, Lernen und schlie\ss{}lich Bewusstsein nennen \citep{FristonFreeEnergy}.

% ===================================================================
\section{Testbare Vorhersagen}
\label{sec:predictions}
% ===================================================================

\subsection{P1: Proteinfaltungslandschaften}

\begin{proposition}
Die Energielandschaft der Proteinfaltung kann durch Nash-Gleichgewichts-Stabilit\"atsanalyse charakterisiert werden.
\end{proposition}

\textbf{Test:} Wende spieltheoretische Stabilit\"atsanalyse (Hesse-Eigenwerte) auf bekannte Proteinfaltungslandschaften an. Sage Faltungsintermediate und Fehlfaltungsfallen vorher.

\subsection{P2: Genregulationsnetzwerk-Attraktoren}

\begin{proposition}
Zelltyp-Attraktoren in Genregulationsnetzwerken k\"onnen aus der Nash-Analyse des TF-DNA-Bindungsspiels vorhergesagt werden.
\end{proposition}

\textbf{Test:} Berechne f\"ur gut charakterisierte GRNs Nash-Gleichgewichte. Vergleiche vorhergesagte Attraktorzust\"ande mit beobachteten Zelltypen.

\subsection{P3: Krebs und Koalitionsstabilit\"at}

\begin{proposition}
Die Krebsinzidenz korreliert invers mit der multizellul\"aren Koalitionsstabilit\"at.
\end{proposition}

\textbf{Test:} \"Uber Gewebetypen hinweg, quantifiziere die Koalitionsstabilit\"at mittels spieltheoretischer Metriken. Sage Krebsinzidenzraten vorher.

\subsection{P4: Metabolische Rate und Fortpflanzungserfolg}

\begin{proposition}
\"Uber Arten hinweg korreliert die massenspezifische Stoffwechselrate mit der Fortpflanzungsrate nach Kontrolle f\"ur K\"orpergr\"o\ss{}e.
\end{proposition}

\textbf{Test:} Stelle Stoffwechselraten und Fortpflanzungsleistungen \"uber Taxa zusammen. Eine positive Korrelation unterst\"utzt die MEPP-Fitness-Verbindung.

\subsection{P5: \"Okosystem-Resilienz}

\begin{proposition}
\"Okosystem-Resilienz kann aus der Nash-Stabilit\"at des \"okologischen Gleichgewichts vorhergesagt werden.
\end{proposition}

\textbf{Test:} Berechne Nash-Stabilit\"at f\"ur gut charakterisierte \"Okosysteme. Vergleiche mit empirischen Resilienzma\ss{}en.

% ===================================================================
\section{Diskussion}
\label{sec:discussion}
% ===================================================================

\subsection{Was ist wirklich neu?}

Die klassische evolution\"are Spieltheorie ist auf der organismischen Ebene gut etabliert. Die Beitr\"age dieses Papers sind:
\begin{enumerate}
    \item \textbf{Vertikale Integration:} Verbindung molekularer, zellul\"arer, organismischer und \"Okosystem-Spiele
    \item \textbf{Thermodynamische Fundierung:} Verankerung von Fitness in MEPP und dissipativer Adaptation
    \item \textbf{RG-Interpretation:} Formalisierung von Haupt\"uberg\"angen als Grobk\"ornung
    \item \textbf{MEPP als Stabilit\"atsfilter:} Reformulierung von MEPP als Stabilit\"atskriterium statt Maximierungsprinzip
\end{enumerate}

\subsection{Unterbrochenes Gleichgewicht als Resolventen\"ubergang}

\begin{remark}[Biologische Stasis und pl\"otzliche \"Uberg\"ange]
\label{rem:punctuated-equilibrium}
Das Muster langer biologischer Stasis, unterbrochen von pl\"otzlichen evolution\"aren \"Uberg\"angen (Punctuated Equilibrium, \citealp{KauffmanOrigins}), findet eine strukturelle Erkl\"arung im Resolventenrahmen \citep{FST-RH2}. Perioden der Stasis entsprechen resolventenstabilen Phasen: die Kopplungsstruktur zweiter Ordnung erh\"alt die aktuelle Konfiguration gegen St\"orungen. Pl\"otzliche \"Uberg\"ange treten auf, wenn Umwelt- oder interne \"Anderungen die Kopplungsarchitektur zweiter Ordnung reorganisieren, analog zum Phasen\"ubergang, der in der Resolventenl\"ucke bei kritischen Parametern in der Riemannschen-Vermutungs-Analyse beobachtet wird (wo $\rho(\lambda)$ Null durchquert und eine qualitative \"Anderung im Stabilit\"atsmechanismus ausl\"ost). Evolution\"are Haupt\"uberg\"ange---vom Prokaryoten zum Eukaryoten, von der Einzelzelligkeit zur Multizellularit\"at---sind somit keine graduellen Optimierungen, sondern \emph{Resolventen\"uberg\"ange}: Reorganisationen der Kopplungsstruktur, die einen qualitativ neuen stabilen Fixpunkt etablieren.
\end{remark}

\subsection{Stabilit\"at vs.\ Maximierung: Die Resolventenperspektive}

Das FST-Framework beschreibt \emph{Stabilit\"atsprinzipien}, keine Maximierungsprinzipien. Die Dynamik f\"uhrender Ordnung ist entartet: viele biologische Konfigurationen sind in erster Ordnung gleicherma\ss{}en lebensfähig. Die Selektion erfolgt in zweiter Ordnung \"uber resolventenartige Kopplung zwischen entarteten und komplement\"aren Unterr\"aumen. MEPP wirkt als Stabilit\"atsfilter, der den einzigen resolventenstabilen Fixpunkt innerhalb der entarteten Mannigfaltigkeit identifiziert. Diese strukturelle Logik---motiviert durch die strukturellen Parallelen in der eichtheoretischen Reformulierung der Riemannschen Vermutung \citep{FST-RH2}---adressiert gleichzeitig die MEPP-Kontroverse (MEPP selektiert f\"ur Stabilit\"at, nicht f\"ur maximale Dissipation), die ESS-Interpretation (ESS sind resolventenstabile Fixpunkte, keine Fitnessoptima), die Rolle der Nash-Frustration (ein notwendiges strukturelles Merkmal der Stabilisierung zweiter Ordnung) und das unterbrochene Gleichgewicht (pl\"otzliche \"Uberg\"ange reflektieren Reorganisation der Resolventen-Kopplungsstruktur). \"Uber alle drei FST-Papers hinweg gilt die gleiche Selektionslogik zweiter Ordnung: Nash-Gleichgewichte sind keine Optima, sondern resolventenstabile Fixpunkte.

\subsection{MEPP: Geltungsbereich und Einschr\"ankungen}

Das Maximum-Entropieproduktions-Prinzip wird in diesem Paper als heuristisches Organisationsprinzip verwendet, nicht als bewiesenes physikalisches Gesetz. Wichtige Einschr\"ankungen m\"ussen anerkannt werden:
\begin{itemize}
\item Die theoretischen Grundlagen von MEPP erfordern lokales thermodynamisches Gleichgewicht und bilineare Entropieproduktion \citep{Martyushev2010}. Biologische Systeme fern vom Gleichgewicht verletzen diese Bedingungen in der Regel.
\item Die Resolvent-Perspektive (Abschnitt~\ref{sec:discussion}) adressiert diese Einschr\"ankung, indem MEPP als Stabilit\"atsfilter reinterpretiert wird: Unter den vielen biologisch lebensf\"ahigen Konfigurationen erster Ordnung identifiziert MEPP den resolvent-stabilen Fixpunkt -- nicht die Konfiguration mit maximaler Dissipation, sondern diejenige, bei der alle Destabilisierungskan\"ale zweiter Ordnung gleichzeitig unterdr\"uckt sind.
\item Diese schw\"achere Interpretation ist kompatibel mit Prigogines minimaler Entropieproduktion nahe dem Gleichgewicht: Intern k\"onnen Organismen die Entropieproduktion minimieren (FEP), w\"ahrend sie extern maximieren (MEPP), wie in Proposition~\ref{prop:fep-mepp} gezeigt.
\end{itemize}

\paragraph{Zirkularit\"atsvorbehalt.} Wenn MEPP als Payoff-Funktion dient und Nash-Gleichgewichte als MEPP-maximierende Zust\"ande definiert werden, ist die Behauptung, dass biologische Systeme Nash-Gleichgewichte einnehmen, \emph{weil} sie Entropieproduktion maximieren, definitorisch zirkul\"ar. Der nicht-zirkul\"are Gehalt dieses Frameworks ist die \emph{Stabilit\"ats}-Aussage: Biologische Attraktoren sind nicht blo\ss{} entropieproduzierend, sondern in zweiter Ordnung stabile Fixpunkte, bei denen alle Destabilisierungskan\"ale (Parasitismus, Mutation, \"okologische Invasion) gleichzeitig unterdr\"uckt sind. Diese Stabilit\"atseigenschaft ist unabh\"angig von der Payoff-Definition testbar (Vorhersagen P1--P5).

\subsection{Ist Biologie ``nur Physik''?}

Die Antwort ist differenziert:
\begin{itemize}
    \item \textbf{Ja:} Das gleiche Optimierungsprinzip (MEPP) und die gleiche mathematische Struktur (Nash-Gleichgewicht) gelten \"uber Skalen hinweg
    \item \textbf{Nein:} Emergenz ist real---h\"ohere Ebenen haben genuine kausale M\"achte, die nicht allein aus niedrigeren Ebenen vorhersagbar sind
\end{itemize}

Der Rahmen ist nicht reduktionistisch im eliminativistischen Sinne. Er identifiziert eine gemeinsame logische Struktur, die auf jeder Skala unterschiedlich instanziiert wird.

% ===================================================================
\section{Schlussfolgerung}
\label{sec:conclusion}
% ===================================================================

FST-III zeigt, dass die Trennung zwischen harten physikalischen Gesetzen und den ``weichen'' Strategien des Lebens k\"unstlich ist. Biologie ist angewandte Thermodynamik und Spieltheorie in einem kosmologisch abgestimmten Rahmen.

Das Maximum Entropy Production Principle liefert den Antrieb, die dissipative Adaptation den Mechanismus der Formbildung, die spieltheoretische ESS die Kriterien f\"ur Langzeitstabilit\"at und das Free Energy Principle die Grundlage f\"ur adaptives Verhalten.

Die prim\"aren Beitr\"age von FST-III sind biologisch: der Nash-Frustrations-Formalismus f\"ur Proteinstrukturen, die hierarchische Spielinterpretation evolution\"arer Haupt\"uberg\"ange und das Koalitionszusammenbruch-Modell der Krebsprogression. Diese konstituieren einen koh\"arenten und testbaren Rahmen.

Die spekulative Erweiterung auf kosmologische Skalen (Abschnitt~\ref{sec:cosmology}) verbindet sich mit dem umfassenderen FST-Forschungsprogramm, erfordert aber dedizierte formale Entwicklung, bevor sie als Teil des etablierten Rahmens betrachtet werden kann.

F\"ur zuk\"unftige Forschung sind die unmittelbaren Priorit\"aten: (i) empirische Validierung der Nash-Frustrationswerte gegen allosterische Datens\"atze; (ii) Kalibrierung des Lernraten-Parameters $\eta$; (iii) Anwendung der Koalitionsstabilit\"atsmetrik (Proposition~\ref{prop:multicellularity}) auf Krebsinzidenzdaten; und (iv) Formalisierung der RG-artigen Grobk\"ornungsschritte als explizite Renormierungsfl\"usse, sofern die Mathematik dies erlaubt.

Wenn die kosmologische Konjektur schlie\ss{}lich untermauert werden kann, w\"urde sie biologische Ph\"anomene als Teil einer kontinuierlichen Kette von Optimierungsprozessen einordnen, die die Geschichte des Universums umspannen. Von der Aminos\"aure bis zum \"Okosystem h\"alt die spieltheoretische Logik. Ob die gleiche Logik bis zu galaktischen Skalen reicht, bleibt eine offene Frage.

% ===================================================================
% Literaturverzeichnis
% ===================================================================

\bibliographystyle{plainnat}
\begin{thebibliography}{99}

\bibitem[Geiger(2026a)]{FSTI}
Geiger, L. (2026).
\newblock Thermodynamic game theory of fundamental particles: The Standard Model as Nash-optimal equilibrium.
\newblock Preprint (FST-I).

\bibitem[Geiger(2026b)]{FSTII}
Geiger, L. (2026).
\newblock Chemical game theory: Prebiotic chemistry as thermodynamic Nash equilibrium under MEPP.
\newblock Preprint (FST-II).

\bibitem[Aktipis et al.(2015)]{AktipisCancer}
Aktipis, C.~A. et al. (2015).
\newblock Cancer across the tree of life.
\newblock \textit{Phil. Trans. R. Soc. B}, 370, 20140219.

\bibitem[Anfinsen(1973)]{Anfinsen1973}
Anfinsen, C.~B. (1973).
\newblock Principles that govern the folding of protein chains.
\newblock \textit{Science}, 181(4096), 223--230.

\bibitem[Graudenzi et al.(2011)]{BooleanEvolution}
Graudenzi, A., Serra, R., Villani, M., Damiani, C., Colacci, A. \& Kauffman, S.~A. (2011).
\newblock Dynamical properties of a Boolean model of gene regulatory network with memory.
\newblock \textit{Journal of Computational Biology}, 18(10), 1291--1303.

\bibitem[M\"ussel et al.(2010)]{BooleanMicroarray}
M\"ussel, C., Hopfensitz, M. \& Kestler, H.~A. (2010).
\newblock BoolNet---an R package for generation, reconstruction and analysis of Boolean networks.
\newblock \textit{Bioinformatics}, 26(10), 1378--1380.

\bibitem[Kauffman(1969)]{BooleanPost}
Kauffman, S.~A. (1969).
\newblock Metabolic stability and epigenesis in randomly constructed genetic nets.
\newblock \textit{Journal of Theoretical Biology}, 22(3), 437--467.

\bibitem[Kauffman(1993)]{KauffmanOrigins}
Kauffman, S.~A. (1993).
\newblock \textit{The Origins of Order: Self-Organization and Selection in Evolution}.
\newblock Oxford University Press.

\bibitem[Hofbauer \& Sigmund(1998)]{CoexistenceLotkaVolterra}
Hofbauer, J. \& Sigmund, K. (1998).
\newblock \textit{Evolutionary Games and Population Dynamics}.
\newblock Cambridge University Press. (Ch. 7: Lotka-Volterra and replicator equivalence).

\bibitem[Milgrom(1983)]{DarkMatterMOND}
Milgrom, M. (1983).
\newblock A modification of the Newtonian dynamics as a possible alternative to the hidden mass hypothesis.
\newblock \textit{The Astrophysical Journal}, 270, 365--370.
% Note: The Hunt Thesis conjecture discussed in Section 8 is the author's own unpublished hypothesis; Milgrom (1983) provides the foundational MOND reference only.

\bibitem[England(2013)]{EnglandJCP}
England, J.~L. (2013).
\newblock Statistical physics of self-replication.
\newblock \textit{Journal of Chemical Physics}, 139, 121923.

\bibitem[England(2015)]{EnglandNano}
England, J.~L. (2015).
\newblock Dissipative adaptation in driven self-assembly.
\newblock \textit{Nature Nanotechnology}, 10, 919--923.

\bibitem[Kleidon(2010)]{EcosystemMEPP}
Kleidon, A. (2010).
\newblock Life, hierarchy, and the thermodynamic machinery of planet Earth.
\newblock \textit{Physics of Life Reviews}, 7(4), 424--460.

\bibitem[Bryngelson \& Wolynes(1989)]{EnzymeFolding}
Bryngelson, J.~D. \& Wolynes, P.~G. (1989).
\newblock Intermediates and barrier crossing in a random energy model (with applications to protein folding).
\newblock \textit{The Journal of Physical Chemistry}, 93(19), 6902--6915.

\bibitem[Friston(2009)]{FristonFreeEnergy}
Friston, K. (2009).
\newblock The free-energy principle: a rough guide to the brain?
\newblock \textit{Trends in Cognitive Sciences}, 13(7), 293--301.

\bibitem[Friston(2010)]{FristonPDF}
Friston, K. (2010).
\newblock The free-energy principle: a unified brain theory?
\newblock \textit{Nature Reviews Neuroscience}, 11, 127--138.

\bibitem[Hamilton(1964)]{Hamilton1964}
Hamilton, W.~D. (1964).
\newblock The genetical evolution of social behaviour.
\newblock \textit{Journal of Theoretical Biology}, 7(1), 1--16.

\bibitem[Maynard Smith \& Price(1973)]{MaynardSmithPrice1973}
Maynard Smith, J. \& Price, G.~R. (1973).
\newblock The logic of animal conflict.
\newblock \textit{Nature}, 246, 15--18.

\bibitem[Bohl et al.(2004)]{GameBiochemistry}
Bohl, K., Hummert, S., Werner, S., Bergfeld, D., Ziesche, A., Fischer, J., Aravind, L. \& Obsil, T. (2004).
\newblock Evolutionary game theory: molecules as players.
\newblock \textit{Molecular BioSystems}, 10(12), 3066--3074.

\bibitem[Maynard Smith(1982)]{MaynardSmith1982}
Maynard Smith, J.
\newblock \textit{Evolution and the Theory of Games}.
\newblock Cambridge UP, 1982.

\bibitem[Holt \& Roth(2004)]{NashPNAS}
Holt, C.~A. \& Roth, A.~E. (2004).
\newblock The Nash equilibrium: A perspective.
\newblock \textit{PNAS}, 101(12), 3999--4002.

\bibitem[Nash(1950)]{NashSFI}
Nash, J.~F. (1950).
\newblock Equilibrium points in n-person games.
\newblock \textit{PNAS}, 36(1), 48--49.

\bibitem[Cressman \& Garay(2003)]{PredatorPreyESS}
Cressman, R. \& Garay, J. (2003).
\newblock Evolutionary stability in Lotka-Volterra systems.
\newblock \textit{Journal of Theoretical Biology}, 222(2), 233--245.

\bibitem[McGaugh et al.(2016)]{RARGalaxies}
McGaugh, S.~S., Lelli, F. \& Schombert, J.~M. (2016).
\newblock Radial acceleration relation in rotationally supported galaxies.
\newblock \textit{Physical Review Letters}, 117(20), 201101.

\bibitem[Sandhu et al.(2026)]{ReplicatorOptimization}
Sandhu, R. et al. (2026).
\newblock The replicator-optimization mechanism: A scale-relative formalism.
\newblock arXiv:2601.06363.

\bibitem[Safron et al.(2025)]{ScaleRGFrontiers}
Safron, A. et al. (2025).
\newblock From pixels to planning: scale-free active inference.
\newblock \textit{Front. Netw. Physiol.}.

\bibitem[Wilson(1975)]{ScaleRGResearch}
Wilson, K.~G. (1975).
\newblock The renormalization group: critical phenomena and the Kondo problem.
\newblock \textit{Reviews of Modern Physics}, 47(4), 773--840.

\bibitem[Schr\"odinger(1944)]{SchoolBiologicalPhysics}
Schr\"odinger, E. (1944).
\newblock \textit{What is Life? The Physical Aspect of the Living Cell}.
\newblock Cambridge University Press.

\bibitem[Gallos et al.(2012)]{SelfSimilarNetworks}
Gallos, L.~K., Makse, H.~A. \& Sigman, M. (2012).
\newblock A small world of weak ties provides optimal global integration of self-similar modules in functional brain networks.
\newblock \textit{J. R. Soc. Interface}, 9, 2131.

\bibitem[Szathmary \& Maynard Smith(1995)]{SzathmaryMaynardSmith1995}
Szathmary, E. \& Maynard Smith, J.
\newblock The major evolutionary transitions.
\newblock \textit{Nature}, 374, 227, 1995.

\bibitem[Ferreiro(2007)]{Ferreiro2007}
Ferreiro, D.U., Hegler, J.A., Komives, E.A. \& Wolynes, P.G.
\newblock Localizing frustration in native proteins and protein assemblies.
\newblock \textit{PNAS}, 104, 19819--19824, 2007.

\bibitem[Ferreiro(2014)]{Ferreiro2014}
Ferreiro, D.U., Komives, E.A. \& Wolynes, P.G.
\newblock Frustration in biomolecules.
\newblock \textit{Quarterly Reviews of Biophysics}, 47, 285--363, 2014.

\bibitem[Martyushev(2010)]{Martyushev2010}
Martyushev, L.~M. (2010).
\newblock The maximum entropy production principle: two basic questions.
\newblock \textit{Phil.\ Trans.\ R.\ Soc.\ B}, 365, 1333--1334.

\bibitem[Geiger(2026c)]{FST-RH3}
Geiger, L. (2026).
\newblock The Riemann Hypothesis Part III: The Analytical Rank-Finite Certificate.
\newblock \textit{Zenodo preprint}, March 2026.

\bibitem[Geiger(2026e)]{FST-RH2}
Geiger, L. (2026).
\newblock Even Dominance of the Weil Quadratic Form: Spectral Analysis, Computer-Assisted Proof, and the Resolvent Gap Theorem.
\newblock \textit{Zenodo preprint}, March 2026.

\bibitem[Geiger(2026d)]{FST-Unified2026}
Geiger, L. (2026).
\newblock The Renormalized Free Energy Framework: A Unified Resolution of Structural Bottlenecks.
\newblock \textit{FST Framework Paper}, 2026.

\end{thebibliography}

\section*{KI-Offenlegung}
Diese Arbeit wurde von Claude (Anthropic) bei der Literatursynthese, mathematischen Verifikation und Texterstellung unterst\"utzt. Alle wissenschaftlichen Behauptungen, originalen Argumente und theoretischen Beitr\"age liegen in der alleinigen Verantwortung des menschlichen Autors. Kein KI-System ist als Koautor gelistet.

\end{document}
